\section{Test Case 23: Refinement with Collector Removal}

Test Case 23 extends the repeated-refinement path through the collector. It
retains Maxwell rheology, Yarin evaporation, stochastic Platen aerodynamic
drag, nozzle insertion, and ordinary numerical anchor beads. The initial jet
contains 400 elements over the first 8 cm of a 12 cm nozzle-to-collector
domain. The applied potential is scaled to retain the physical field strength
of Test Case 22. The timestep is $2.5\times10^{-8}$ s and the run ends after
16,000 steps.

The dedicated runner sets developer-only initial and growth reserves of 20
entries. Three accepted Akima events therefore cause three capacity changes.
Collector removal changes the lower active bound before the third remesh and
continues afterwards. This combines all relevant topology transitions:
insertion, removal, remeshing, capacity replacement, Platen workspace resize,
and semantic repacking of the pre-generated Gaussian history.

With NVFORTRAN 24.3, the CPU refinement events occur at steps 14,301, 14,944,
and 15,718. The first collector removal occurs at step 15,647. The complete
force oracle places the last remesh at step 15,717, while the native NVIDIA
A30 places it at step 15,700. All paths preserve 81 active anchors and grow
capacity from 420 to 477, then 520, and finally 557.

The exact later removal sequence is not a CPU/GPU identity requirement.
Target-centric direct Coulomb accumulation changes the floating-point order
of summation, and the stochastic bending trajectory amplifies this difference
near the collector. Validation instead requires collector removal both before
and after the final remesh, finite output, ordered coordinates, unchanged
anchor fields, and conservation of reference volume, evaporated volume, mass,
and charge at every Akima event.

Removal is performed by the device topology primitive and does not download
the full jet. Complete state transfers remain limited to accepted refinement
events and final shutdown. Host code constructs the target mesh and applies
conservation and anchor rules; Akima coefficient construction and all 11 field
interpolations execute on the GPU.

An NVIDIA A30 transfer audit records four complete state downloads, namely
the three accepted Akima events and final shutdown, together with exactly
three Gaussian-history uploads, topology rebinds, and evaporation-state
rebinds. The ordinary collector check returns one four-byte control value;
an accepted removal transfers and clears only the collected point.

A dedicated A/B build evaluates the historical host Akima routines and the
device implementation at all three events. Across 33 field/event comparisons,
the maximum coefficient relative difference is $2.48\times10^{-15}$, while
interpolated values differ by at most $7.28\times10^{-12}$ absolutely and
$3.32\times10^{-15}$ relatively. Independent device and host-Akima-oracle
runs produce byte-identical statistics files and the same event and removal
topology.

The commands
\texttt{tests/refinement/run\_test23.sh nvfortran},
\texttt{tests/refinement/run\_test23.sh openacc}, and
\texttt{tests/refinement/run\_test23.sh force-oracle} run the CPU, native GPU,
and force-oracle validations. Arguments \texttt{host-akima} and
\texttt{akima-compare} select the historical spline oracle and the direct
host/device spline comparison, respectively.
